% ============================================================
%  macros.tex — VDFI paper-writing shortcuts
% ------------------------------------------------------------
%  Right after \documentclass{...}, add:
%     % ============================================================
%  macros.tex — VDFI paper-writing shortcuts
% ------------------------------------------------------------
%  Right after \documentclass{...}, add:
%     % ============================================================
%  macros.tex — VDFI paper-writing shortcuts
% ------------------------------------------------------------
%  Right after \documentclass{...}, add:
%     % ============================================================
%  macros.tex — VDFI paper-writing shortcuts
% ------------------------------------------------------------
%  Right after \documentclass{...}, add:
%     \input{macros.tex}
%  (Rename to macros.sty and swap in \usepackage{macros} if you
%  prefer treating it as a package — behaves the same either way.)
% ============================================================

\usepackage{xspace}
\usepackage{amsmath}

% ---- optional: robust unit handling, see latex-tips.md ----
% \usepackage{siunitx}
% \sisetup{per-mode=symbol}

% ------------------------------------------------------------
% Ions & emission lines
% ------------------------------------------------------------
% \ion is built into AASTeX (aastex63/aastex7). If you are NOT
% using an AASTeX class, uncomment this fallback definition:
% \providecommand{\ion}[2]{#1$\;$\textsc{\lowercase{#2}}}

\newcommand{\lya}{Ly$\alpha$\xspace}
\newcommand{\lyb}{Ly$\beta$\xspace}
\newcommand{\Ha}{H$\alpha$\xspace}
\newcommand{\Hb}{H$\beta$\xspace}
\newcommand{\Hg}{H$\gamma$\xspace}
\newcommand{\OIII}{[\ion{O}{3}] $\lambda 5007$\xspace}
\newcommand{\OII}{[\ion{O}{2}] $\lambda 3727$\xspace}
\newcommand{\NII}{[\ion{N}{2}]\xspace}
\newcommand{\CIV}{\ion{C}{4} $\lambda 1549$\xspace}
\newcommand{\HI}{\ion{H}{1}\xspace}

% ------------------------------------------------------------
% Angstrom — works outside math mode, spaces itself correctly
% ------------------------------------------------------------
% \AA is LaTeX's built-in Angstrom symbol, but it's a control
% WORD, so TeX silently swallows the space that follows it:
%   "5007\AA emission"   typesets as   "5007Åemission"  (no space!)
% \xspace fixes this by inserting a space unless the next token
% is punctuation or another space, so you never have to remember
% a trailing "\ " or "{}" yourself.
%
% Two versions:
%   \AAns  — raw symbol, NO auto-space. Use when compounding
%            (e.g. building a unit like erg/s/cm^2/Angstrom).
%   \Ang   — auto-spacing version for standalone use in text.
% Named \Ang (capital A) rather than \ang so it doesn't collide
% with siunitx's \ang{} command, which means "degrees".
\newcommand{\AAns}{\AA}
\newcommand{\Ang}{\AAns\xspace}

% ------------------------------------------------------------
% Approximate / order-of-magnitude symbols
% ------------------------------------------------------------
% \ensuremath lets these drop straight into running text with no
% manual $ $, and won't error if you happen to use them inside an
% existing math environment. See latex-tips.md for \sim vs \approx
% vs \simeq conventions.
\newcommand{\sm}{\ensuremath{\sim}\xspace}          % order-of-magnitude / "roughly": z\sm2
\newcommand{\ap}{\ensuremath{\approx}\xspace}        % numeric approx equality: x\ap3.14
\newcommand{\simeqsym}{\ensuremath{\simeq}\xspace}   % "approximately equal to" / similarity

% ------------------------------------------------------------
% Common astro quantities & units
% ------------------------------------------------------------
\newcommand{\Msun}{M$_{\odot}$\xspace}
\newcommand{\Lsun}{L$_{\odot}$\xspace}
\newcommand{\kpc}{kpc\xspace}
\newcommand{\Mpc}{Mpc\xspace}
\newcommand{\kms}{km\,s$^{-1}$\xspace}
\newcommand{\ergs}{erg\,s$^{-1}$\xspace}
\newcommand{\ergscmsq}{erg\,s$^{-1}$\,cm$^{-2}$\xspace}
\newcommand{\ergscmsqA}{erg\,s$^{-1}$\,cm$^{-2}$\,\AAns$^{-1}$\xspace}
\newcommand{\SFR}{SFR\xspace}
\newcommand{\zsys}{$z_{\rm sys}$\xspace}

% ------------------------------------------------------------
% Area / surface-brightness units
% ------------------------------------------------------------
% Macro names can't contain digits (LaTeX control words are
% letters-only) -- "kpc^2" etc. become \kpcsq / \arcsecsq rather
% than \kpc2. (\ergscmsq / \ergscmsqA above are named this way for
% the same reason -- \ergscm2 would NOT be a valid macro name and
% will corrupt the rest of the preamble parse if you try it.)
% Composed units are spelled out in full rather than nesting other
% \xspace macros, so spacing stays predictable.
\newcommand{\kpcsq}{kpc$^{2}$\xspace}
\newcommand{\arcsecsq}{arcsec$^{2}$\xspace}
\newcommand{\Lkpcsq}{erg\,s$^{-1}$\,kpc$^{-2}$\xspace}          % L / kpc^2 -- physical (deprojected) luminosity surface density
\newcommand{\Farcsecsq}{erg\,s$^{-1}$\,cm$^{-2}$\,arcsec$^{-2}$\xspace}  % F / arcsec^2 -- observed surface brightness

% ------------------------------------------------------------
% Usage, once inside \begin{document}:
%
%   The object is bright in \lya, but faint in \Hb and \OIII.
%   The line sits at 5007\Ang, right next to the doublet.
%   The sample spans z\sm2 to z\sm3.
%   The measured flux is 3.2\ap3.4 $\times$ 10$^{-17}$ \ergscmsq.
%   The halo covers $\sm$40 \kpcsq (14 \arcsecsq at this redshift).
%   Central SB: $10^{40}$ \Lkpcsq; observed SB limit: $10^{-18}$ \Farcsecsq.
% ------------------------------------------------------------

%  (Rename to macros.sty and swap in \usepackage{macros} if you
%  prefer treating it as a package — behaves the same either way.)
% ============================================================

\usepackage{xspace}
\usepackage{amsmath}

% ---- optional: robust unit handling, see latex-tips.md ----
% \usepackage{siunitx}
% \sisetup{per-mode=symbol}

% ------------------------------------------------------------
% Ions & emission lines
% ------------------------------------------------------------
% \ion is built into AASTeX (aastex63/aastex7). If you are NOT
% using an AASTeX class, uncomment this fallback definition:
% \providecommand{\ion}[2]{#1$\;$\textsc{\lowercase{#2}}}

\newcommand{\lya}{Ly$\alpha$\xspace}
\newcommand{\lyb}{Ly$\beta$\xspace}
\newcommand{\Ha}{H$\alpha$\xspace}
\newcommand{\Hb}{H$\beta$\xspace}
\newcommand{\Hg}{H$\gamma$\xspace}
\newcommand{\OIII}{[\ion{O}{3}] $\lambda 5007$\xspace}
\newcommand{\OII}{[\ion{O}{2}] $\lambda 3727$\xspace}
\newcommand{\NII}{[\ion{N}{2}]\xspace}
\newcommand{\CIV}{\ion{C}{4} $\lambda 1549$\xspace}
\newcommand{\HI}{\ion{H}{1}\xspace}

% ------------------------------------------------------------
% Angstrom — works outside math mode, spaces itself correctly
% ------------------------------------------------------------
% \AA is LaTeX's built-in Angstrom symbol, but it's a control
% WORD, so TeX silently swallows the space that follows it:
%   "5007\AA emission"   typesets as   "5007Åemission"  (no space!)
% \xspace fixes this by inserting a space unless the next token
% is punctuation or another space, so you never have to remember
% a trailing "\ " or "{}" yourself.
%
% Two versions:
%   \AAns  — raw symbol, NO auto-space. Use when compounding
%            (e.g. building a unit like erg/s/cm^2/Angstrom).
%   \Ang   — auto-spacing version for standalone use in text.
% Named \Ang (capital A) rather than \ang so it doesn't collide
% with siunitx's \ang{} command, which means "degrees".
\newcommand{\AAns}{\AA}
\newcommand{\Ang}{\AAns\xspace}

% ------------------------------------------------------------
% Approximate / order-of-magnitude symbols
% ------------------------------------------------------------
% \ensuremath lets these drop straight into running text with no
% manual $ $, and won't error if you happen to use them inside an
% existing math environment. See latex-tips.md for \sim vs \approx
% vs \simeq conventions.
\newcommand{\sm}{\ensuremath{\sim}\xspace}          % order-of-magnitude / "roughly": z\sm2
\newcommand{\ap}{\ensuremath{\approx}\xspace}        % numeric approx equality: x\ap3.14
\newcommand{\simeqsym}{\ensuremath{\simeq}\xspace}   % "approximately equal to" / similarity

% ------------------------------------------------------------
% Common astro quantities & units
% ------------------------------------------------------------
\newcommand{\Msun}{M$_{\odot}$\xspace}
\newcommand{\Lsun}{L$_{\odot}$\xspace}
\newcommand{\kpc}{kpc\xspace}
\newcommand{\Mpc}{Mpc\xspace}
\newcommand{\kms}{km\,s$^{-1}$\xspace}
\newcommand{\ergs}{erg\,s$^{-1}$\xspace}
\newcommand{\ergscmsq}{erg\,s$^{-1}$\,cm$^{-2}$\xspace}
\newcommand{\ergscmsqA}{erg\,s$^{-1}$\,cm$^{-2}$\,\AAns$^{-1}$\xspace}
\newcommand{\SFR}{SFR\xspace}
\newcommand{\zsys}{$z_{\rm sys}$\xspace}

% ------------------------------------------------------------
% Area / surface-brightness units
% ------------------------------------------------------------
% Macro names can't contain digits (LaTeX control words are
% letters-only) -- "kpc^2" etc. become \kpcsq / \arcsecsq rather
% than \kpc2. (\ergscmsq / \ergscmsqA above are named this way for
% the same reason -- \ergscm2 would NOT be a valid macro name and
% will corrupt the rest of the preamble parse if you try it.)
% Composed units are spelled out in full rather than nesting other
% \xspace macros, so spacing stays predictable.
\newcommand{\kpcsq}{kpc$^{2}$\xspace}
\newcommand{\arcsecsq}{arcsec$^{2}$\xspace}
\newcommand{\Lkpcsq}{erg\,s$^{-1}$\,kpc$^{-2}$\xspace}          % L / kpc^2 -- physical (deprojected) luminosity surface density
\newcommand{\Farcsecsq}{erg\,s$^{-1}$\,cm$^{-2}$\,arcsec$^{-2}$\xspace}  % F / arcsec^2 -- observed surface brightness

% ------------------------------------------------------------
% Usage, once inside \begin{document}:
%
%   The object is bright in \lya, but faint in \Hb and \OIII.
%   The line sits at 5007\Ang, right next to the doublet.
%   The sample spans z\sm2 to z\sm3.
%   The measured flux is 3.2\ap3.4 $\times$ 10$^{-17}$ \ergscmsq.
%   The halo covers $\sm$40 \kpcsq (14 \arcsecsq at this redshift).
%   Central SB: $10^{40}$ \Lkpcsq; observed SB limit: $10^{-18}$ \Farcsecsq.
% ------------------------------------------------------------

%  (Rename to macros.sty and swap in \usepackage{macros} if you
%  prefer treating it as a package — behaves the same either way.)
% ============================================================

\usepackage{xspace}
\usepackage{amsmath}

% ---- optional: robust unit handling, see latex-tips.md ----
% \usepackage{siunitx}
% \sisetup{per-mode=symbol}

% ------------------------------------------------------------
% Ions & emission lines
% ------------------------------------------------------------
% \ion is built into AASTeX (aastex63/aastex7). If you are NOT
% using an AASTeX class, uncomment this fallback definition:
% \providecommand{\ion}[2]{#1$\;$\textsc{\lowercase{#2}}}

\newcommand{\lya}{Ly$\alpha$\xspace}
\newcommand{\lyb}{Ly$\beta$\xspace}
\newcommand{\Ha}{H$\alpha$\xspace}
\newcommand{\Hb}{H$\beta$\xspace}
\newcommand{\Hg}{H$\gamma$\xspace}
\newcommand{\OIII}{[\ion{O}{3}] $\lambda 5007$\xspace}
\newcommand{\OII}{[\ion{O}{2}] $\lambda 3727$\xspace}
\newcommand{\NII}{[\ion{N}{2}]\xspace}
\newcommand{\CIV}{\ion{C}{4} $\lambda 1549$\xspace}
\newcommand{\HI}{\ion{H}{1}\xspace}

% ------------------------------------------------------------
% Angstrom — works outside math mode, spaces itself correctly
% ------------------------------------------------------------
% \AA is LaTeX's built-in Angstrom symbol, but it's a control
% WORD, so TeX silently swallows the space that follows it:
%   "5007\AA emission"   typesets as   "5007Åemission"  (no space!)
% \xspace fixes this by inserting a space unless the next token
% is punctuation or another space, so you never have to remember
% a trailing "\ " or "{}" yourself.
%
% Two versions:
%   \AAns  — raw symbol, NO auto-space. Use when compounding
%            (e.g. building a unit like erg/s/cm^2/Angstrom).
%   \Ang   — auto-spacing version for standalone use in text.
% Named \Ang (capital A) rather than \ang so it doesn't collide
% with siunitx's \ang{} command, which means "degrees".
\newcommand{\AAns}{\AA}
\newcommand{\Ang}{\AAns\xspace}

% ------------------------------------------------------------
% Approximate / order-of-magnitude symbols
% ------------------------------------------------------------
% \ensuremath lets these drop straight into running text with no
% manual $ $, and won't error if you happen to use them inside an
% existing math environment. See latex-tips.md for \sim vs \approx
% vs \simeq conventions.
\newcommand{\sm}{\ensuremath{\sim}\xspace}          % order-of-magnitude / "roughly": z\sm2
\newcommand{\ap}{\ensuremath{\approx}\xspace}        % numeric approx equality: x\ap3.14
\newcommand{\simeqsym}{\ensuremath{\simeq}\xspace}   % "approximately equal to" / similarity

% ------------------------------------------------------------
% Common astro quantities & units
% ------------------------------------------------------------
\newcommand{\Msun}{M$_{\odot}$\xspace}
\newcommand{\Lsun}{L$_{\odot}$\xspace}
\newcommand{\kpc}{kpc\xspace}
\newcommand{\Mpc}{Mpc\xspace}
\newcommand{\kms}{km\,s$^{-1}$\xspace}
\newcommand{\ergs}{erg\,s$^{-1}$\xspace}
\newcommand{\ergscmsq}{erg\,s$^{-1}$\,cm$^{-2}$\xspace}
\newcommand{\ergscmsqA}{erg\,s$^{-1}$\,cm$^{-2}$\,\AAns$^{-1}$\xspace}
\newcommand{\SFR}{SFR\xspace}
\newcommand{\zsys}{$z_{\rm sys}$\xspace}

% ------------------------------------------------------------
% Area / surface-brightness units
% ------------------------------------------------------------
% Macro names can't contain digits (LaTeX control words are
% letters-only) -- "kpc^2" etc. become \kpcsq / \arcsecsq rather
% than \kpc2. (\ergscmsq / \ergscmsqA above are named this way for
% the same reason -- \ergscm2 would NOT be a valid macro name and
% will corrupt the rest of the preamble parse if you try it.)
% Composed units are spelled out in full rather than nesting other
% \xspace macros, so spacing stays predictable.
\newcommand{\kpcsq}{kpc$^{2}$\xspace}
\newcommand{\arcsecsq}{arcsec$^{2}$\xspace}
\newcommand{\Lkpcsq}{erg\,s$^{-1}$\,kpc$^{-2}$\xspace}          % L / kpc^2 -- physical (deprojected) luminosity surface density
\newcommand{\Farcsecsq}{erg\,s$^{-1}$\,cm$^{-2}$\,arcsec$^{-2}$\xspace}  % F / arcsec^2 -- observed surface brightness

% ------------------------------------------------------------
% Usage, once inside \begin{document}:
%
%   The object is bright in \lya, but faint in \Hb and \OIII.
%   The line sits at 5007\Ang, right next to the doublet.
%   The sample spans z\sm2 to z\sm3.
%   The measured flux is 3.2\ap3.4 $\times$ 10$^{-17}$ \ergscmsq.
%   The halo covers $\sm$40 \kpcsq (14 \arcsecsq at this redshift).
%   Central SB: $10^{40}$ \Lkpcsq; observed SB limit: $10^{-18}$ \Farcsecsq.
% ------------------------------------------------------------

%  (Rename to macros.sty and swap in \usepackage{macros} if you
%  prefer treating it as a package — behaves the same either way.)
% ============================================================

\usepackage{xspace}
\usepackage{amsmath}

% ---- optional: robust unit handling, see latex-tips.md ----
% \usepackage{siunitx}
% \sisetup{per-mode=symbol}

% ------------------------------------------------------------
% Ions & emission lines
% ------------------------------------------------------------
% \ion is built into AASTeX (aastex63/aastex7). If you are NOT
% using an AASTeX class, uncomment this fallback definition:
% \providecommand{\ion}[2]{#1$\;$\textsc{\lowercase{#2}}}

\newcommand{\lya}{Ly$\alpha$\xspace}
\newcommand{\lyb}{Ly$\beta$\xspace}
\newcommand{\Ha}{H$\alpha$\xspace}
\newcommand{\Hb}{H$\beta$\xspace}
\newcommand{\Hg}{H$\gamma$\xspace}
\newcommand{\OIII}{[\ion{O}{3}] $\lambda 5007$\xspace}
\newcommand{\OII}{[\ion{O}{2}] $\lambda 3727$\xspace}
\newcommand{\NII}{[\ion{N}{2}]\xspace}
\newcommand{\CIV}{\ion{C}{4} $\lambda 1549$\xspace}
\newcommand{\HI}{\ion{H}{1}\xspace}

% ------------------------------------------------------------
% Angstrom — works outside math mode, spaces itself correctly
% ------------------------------------------------------------
% \AA is LaTeX's built-in Angstrom symbol, but it's a control
% WORD, so TeX silently swallows the space that follows it:
%   "5007\AA emission"   typesets as   "5007Åemission"  (no space!)
% \xspace fixes this by inserting a space unless the next token
% is punctuation or another space, so you never have to remember
% a trailing "\ " or "{}" yourself.
%
% Two versions:
%   \AAns  — raw symbol, NO auto-space. Use when compounding
%            (e.g. building a unit like erg/s/cm^2/Angstrom).
%   \Ang   — auto-spacing version for standalone use in text.
% Named \Ang (capital A) rather than \ang so it doesn't collide
% with siunitx's \ang{} command, which means "degrees".
\newcommand{\AAns}{\AA}
\newcommand{\Ang}{\AAns\xspace}

% ------------------------------------------------------------
% Approximate / order-of-magnitude symbols
% ------------------------------------------------------------
% \ensuremath lets these drop straight into running text with no
% manual $ $, and won't error if you happen to use them inside an
% existing math environment. See latex-tips.md for \sim vs \approx
% vs \simeq conventions.
\newcommand{\sm}{\ensuremath{\sim}\xspace}          % order-of-magnitude / "roughly": z\sm2
\newcommand{\ap}{\ensuremath{\approx}\xspace}        % numeric approx equality: x\ap3.14
\newcommand{\simeqsym}{\ensuremath{\simeq}\xspace}   % "approximately equal to" / similarity

% ------------------------------------------------------------
% Common astro quantities & units
% ------------------------------------------------------------
\newcommand{\Msun}{M$_{\odot}$\xspace}
\newcommand{\Lsun}{L$_{\odot}$\xspace}
\newcommand{\kpc}{kpc\xspace}
\newcommand{\Mpc}{Mpc\xspace}
\newcommand{\kms}{km\,s$^{-1}$\xspace}
\newcommand{\ergs}{erg\,s$^{-1}$\xspace}
\newcommand{\ergscmsq}{erg\,s$^{-1}$\,cm$^{-2}$\xspace}
\newcommand{\ergscmsqA}{erg\,s$^{-1}$\,cm$^{-2}$\,\AAns$^{-1}$\xspace}
\newcommand{\SFR}{SFR\xspace}
\newcommand{\zsys}{$z_{\rm sys}$\xspace}

% ------------------------------------------------------------
% Area / surface-brightness units
% ------------------------------------------------------------
% Macro names can't contain digits (LaTeX control words are
% letters-only) -- "kpc^2" etc. become \kpcsq / \arcsecsq rather
% than \kpc2. (\ergscmsq / \ergscmsqA above are named this way for
% the same reason -- \ergscm2 would NOT be a valid macro name and
% will corrupt the rest of the preamble parse if you try it.)
% Composed units are spelled out in full rather than nesting other
% \xspace macros, so spacing stays predictable.
\newcommand{\kpcsq}{kpc$^{2}$\xspace}
\newcommand{\arcsecsq}{arcsec$^{2}$\xspace}
\newcommand{\Lkpcsq}{erg\,s$^{-1}$\,kpc$^{-2}$\xspace}          % L / kpc^2 -- physical (deprojected) luminosity surface density
\newcommand{\Farcsecsq}{erg\,s$^{-1}$\,cm$^{-2}$\,arcsec$^{-2}$\xspace}  % F / arcsec^2 -- observed surface brightness

% ------------------------------------------------------------
% Usage, once inside \begin{document}:
%
%   The object is bright in \lya, but faint in \Hb and \OIII.
%   The line sits at 5007\Ang, right next to the doublet.
%   The sample spans z\sm2 to z\sm3.
%   The measured flux is 3.2\ap3.4 $\times$ 10$^{-17}$ \ergscmsq.
%   The halo covers $\sm$40 \kpcsq (14 \arcsecsq at this redshift).
%   Central SB: $10^{40}$ \Lkpcsq; observed SB limit: $10^{-18}$ \Farcsecsq.
% ------------------------------------------------------------
